\documentclass[11pt,a4paper]{article}

% ---- Packages ----
\usepackage[utf8]{inputenc}
\usepackage[T1]{fontenc}
\usepackage{amsmath,amssymb,amsthm}
\usepackage{mathtools}
\usepackage[margin=1in,headheight=14pt]{geometry}
\usepackage{hyperref}
\usepackage{enumitem}
\usepackage{xcolor}
\usepackage{tcolorbox}
\usepackage{fancyhdr}
\usepackage{booktabs}

% ---- Hyperref ----
\hypersetup{
    colorlinks=true,
    linkcolor=red,
    citecolor=blue,
    urlcolor=blue
}

% ---- Header ----
\pagestyle{fancy}
\fancyhf{}
\fancyhead[L]{\textit{Lesson 1: Measure Spaces}}
\fancyhead[R]{\thepage}
\renewcommand{\headrulewidth}{0.4pt}

% ---- Theorem Environments ----
\theoremstyle{definition}
\newtheorem{definition}{Definition}[section]
\newtheorem{example}{Example}[section]
\newtheorem{exercise}{Exercise}[section]

\theoremstyle{plain}
\newtheorem{theorem}{Theorem}[section]
\newtheorem{lemma}[theorem]{Lemma}
\newtheorem{proposition}[theorem]{Proposition}
\newtheorem{corollary}[theorem]{Corollary}

\theoremstyle{remark}
\newtheorem{remark}{Remark}[section]

% ---- Colored Boxes ----
\tcbuselibrary{theorems,skins,breakable}

\newtcolorbox{keyresult}[1][]{
    colback=green!5!white,
    colframe=green!50!black,
    fonttitle=\bfseries,
    title=#1,
    breakable
}

\newtcolorbox{rlconnection}[1][]{
    colback=blue!5!white,
    colframe=blue!50!black,
    fonttitle=\bfseries,
    title=#1,
    breakable
}

\newtcolorbox{warning}[1][]{
    colback=red!5!white,
    colframe=red!50!black,
    fonttitle=\bfseries,
    title=#1,
    breakable
}

% ---- Operators ----
\newcommand{\R}{\mathbb{R}}
\newcommand{\N}{\mathbb{N}}
\newcommand{\Q}{\mathbb{Q}}
\newcommand{\Z}{\mathbb{Z}}
\newcommand{\Prob}{\mathbb{P}}
\newcommand{\E}{\mathbb{E}}
\newcommand{\indicator}{\mathbf{1}}
\newcommand{\powerset}{\mathcal{P}}
\newcommand{\Borel}{\mathcal{B}}

% ---- Title ----
\title{Lesson 2: Measures and the Carath\\{e}odory Extension}
\author{Curriculum: A Rigorous Learning Path to Multi-Agent Deep RL}
\date{February 9, 2026}

\begin{document}
\maketitle
\tableofcontents
\newpage

\setcounter{section}{2}

\section{Measures}
\label{sec:measures}

\subsection{Definition and Fundamental Properties}
\label{subsec:measure_def}

\begin{definition}[Measure]
\label{def:measure}
A \emph{measure} on a measurable space $(\Omega, \mathcal{F})$ is a function
$\mu: \mathcal{F} \to [0, \infty]$ satisfying:
\begin{enumerate}[nosep,label=(\roman*)]
    \item $\mu(\emptyset) = 0$.
    \item \textbf{Countable additivity:} If $A_1, A_2, \ldots \in \mathcal{F}$ are
    pairwise disjoint, then:
    \begin{equation}
        \mu\!\left(\bigsqcup_{n=1}^\infty A_n\right) = \sum_{n=1}^\infty \mu(A_n).
        \label{eq:countable_additivity}
    \end{equation}
\end{enumerate}
The triple $(\Omega, \mathcal{F}, \mu)$ is a \emph{measure space}. If
$\mu(\Omega) = 1$, then $\mu$ is a \emph{probability measure} and we write
$\Prob$ instead of $\mu$.
\end{definition}

\begin{example}[Standard Measures]
\label{ex:measures}
\begin{enumerate}[nosep,label=(\alph*)]
    \item \textbf{Counting measure:} $\mu(A) = |A|$ (number of elements) on any countable $\Omega$.
    \item \textbf{Dirac measure:} $\delta_x(A) = \indicator_A(x) = \begin{cases} 1 & \text{if } x \in A \\ 0 & \text{if } x \notin A \end{cases}$.
    \item \textbf{Lebesgue measure:} $\lambda((a,b]) = b - a$ on $(\R, \Borel(\R))$. (Existence proved in Section~\ref{sec:construction}.)
    \item \textbf{Probability on a finite set:} $\Prob(\{s\}) = p(s)$ with $\sum_{s \in \mathcal{S}} p(s) = 1$.
\end{enumerate}
\end{example}

\begin{theorem}[Properties of Measures]
\label{thm:measure_props}
Let $(\Omega, \mathcal{F}, \mu)$ be a measure space. For $A, B, A_n \in \mathcal{F}$:
\begin{enumerate}[nosep,label=(\alph*)]
    \item \textbf{Monotonicity:} $A \subseteq B \implies \mu(A) \leq \mu(B)$.
    \item \textbf{Subtractivity:} $A \subseteq B$ and $\mu(A) < \infty \implies \mu(B \setminus A) = \mu(B) - \mu(A)$.
    \item \textbf{Sub-additivity:} $\mu\!\left(\bigcup_{n=1}^\infty A_n\right) \leq \sum_{n=1}^\infty \mu(A_n)$.
    \item \textbf{Continuity from below:} $A_n \uparrow A \implies \mu(A) = \lim_n \mu(A_n)$.
    \item \textbf{Continuity from above:} $A_n \downarrow A$ and $\mu(A_1) < \infty \implies \mu(A) = \lim_n \mu(A_n)$.
\end{enumerate}
\end{theorem}

\begin{proof}
\textbf{(a)} $B = A \sqcup (B \setminus A)$, so $\mu(B) = \mu(A) + \mu(B \setminus A) \geq \mu(A)$.

\textbf{(b)} From $\mu(B) = \mu(A) + \mu(B \setminus A)$ and $\mu(A) < \infty$, subtract $\mu(A)$.

\textbf{(c)} Define $B_1 = A_1$, $B_n = A_n \setminus \bigcup_{k=1}^{n-1} A_k$ for $n \geq 2$. The $B_n$ are disjoint, $\bigcup B_n = \bigcup A_n$, and $B_n \subseteq A_n$. So:
$\mu(\bigcup A_n) = \mu(\bigsqcup B_n) = \sum \mu(B_n) \leq \sum \mu(A_n)$.

\textbf{(d)} Let $A_0 = \emptyset$. The sets $C_n = A_n \setminus A_{n-1}$ are disjoint and $A_n = \bigsqcup_{k=1}^n C_k$. Also $A = \bigsqcup_{k=1}^\infty C_k$. By countable additivity:
\begin{equation}
    \mu(A) = \sum_{k=1}^\infty \mu(C_k) = \lim_{n \to \infty} \sum_{k=1}^n \mu(C_k) = \lim_{n \to \infty} \mu(A_n).
\end{equation}

\textbf{(e)} Define $B_n = A_1 \setminus A_n$. Then $B_n \uparrow A_1 \setminus A$. By (d):
$\mu(A_1 \setminus A) = \lim \mu(B_n) = \lim (\mu(A_1) - \mu(A_n))$, using (b) (valid since $\mu(A_n) \leq \mu(A_1) < \infty$). So $\mu(A_1) - \mu(A) = \mu(A_1) - \lim \mu(A_n)$, giving $\mu(A) = \lim \mu(A_n)$.
\end{proof}

\begin{warning}[Finiteness Condition in (e)]
The condition $\mu(A_1) < \infty$ in continuity from above is essential. Consider
counting measure on $\N$ with $A_n = \{n, n+1, n+2, \ldots\}$. Then $A_n \downarrow \emptyset$
but $\mu(A_n) = \infty$ for all $n$, while $\mu(\emptyset) = 0$. For probability
measures, this condition is always satisfied since $\Prob(A_1) \leq 1 < \infty$.
\end{warning}

\subsection{Null Sets and Complete Measures}
\label{subsec:null_sets}

\begin{definition}[Null Set, Complete Measure]
\label{def:null_complete}
A set $N \in \mathcal{F}$ is a \emph{$\mu$-null set} (or \emph{null set}) if $\mu(N) = 0$.
A measure $\mu$ is \emph{complete} if whenever $N$ is $\mu$-null and $A \subseteq N$,
then $A \in \mathcal{F}$ (and hence $\mu(A) = 0$).
\end{definition}

\begin{proposition}[Completion of a Measure]
\label{prop:completion}
Every measure space $(\Omega, \mathcal{F}, \mu)$ can be completed: there exists
a smallest $\sigma$-algebra $\overline{\mathcal{F}} \supseteq \mathcal{F}$ and
a measure $\bar{\mu}$ extending $\mu$ such that $(\Omega, \overline{\mathcal{F}}, \bar{\mu})$
is complete. Explicitly:
\begin{equation}
    \overline{\mathcal{F}} = \{A \cup N' : A \in \mathcal{F},\; N' \subseteq N \text{ for some $\mu$-null } N \in \mathcal{F}\},
\end{equation}
and $\bar{\mu}(A \cup N') = \mu(A)$.
\end{proposition}

\begin{proof}
We verify that $\overline{\mathcal{F}}$ is a $\sigma$-algebra. The key step is
closure under complement: if $E = A \cup N'$ with $N' \subseteq N$ and $\mu(N) = 0$,
then $E^c = A^c \cap (N')^c = (A^c \cap N^c) \cup (A^c \cap N \cap (N')^c)$.
The first part is in $\mathcal{F}$, and the second is a subset of $N$ (a null set),
so $E^c \in \overline{\mathcal{F}}$.

For well-definedness of $\bar{\mu}$: if $A_1 \cup N_1' = A_2 \cup N_2'$, then
$A_1 \triangle A_2 \subseteq N_1 \cup N_2$ (a null set), so
$\mu(A_1 \triangle A_2) = 0$, hence $\mu(A_1) = \mu(A_2)$.
\end{proof}

%% ============================================================
%% SECTION 4: CONSTRUCTION OF MEASURES
%% ============================================================

\section{Construction of Measures: Carath\'{e}odory's Theorem}
\label{sec:construction}

How do we actually \emph{build} measures? We cannot simply assign
$\mu(A)$ for every $A \in \Borel(\R)$ individually --- the Borel $\sigma$-algebra
is uncountably large. Instead, we specify $\mu$ on a small, manageable collection
and extend it.

\subsection{Algebras and Pre-Measures}
\label{subsec:algebras}

\begin{definition}[Algebra of Sets]
\label{def:algebra}
A collection $\mathcal{A} \subseteq \powerset(\Omega)$ is an \emph{algebra}
(or \emph{field}) if:
\begin{enumerate}[nosep,label=(\roman*)]
    \item $\Omega \in \mathcal{A}$.
    \item $A \in \mathcal{A} \implies A^c \in \mathcal{A}$.
    \item $A, B \in \mathcal{A} \implies A \cup B \in \mathcal{A}$.
    \hfill(closed under \emph{finite} union)
\end{enumerate}
\end{definition}

\begin{remark}
An algebra is like a $\sigma$-algebra but with only finite (not countable) closure.
Every $\sigma$-algebra is an algebra, but not conversely.
\end{remark}

\begin{example}
\label{ex:algebra}
On $\R$, the collection $\mathcal{A}$ of all finite unions of half-open intervals
$(a, b]$ (including $\emptyset$ and $\R = (-\infty, \infty)$) is an algebra
but not a $\sigma$-algebra. For instance, $\bigcup_{n=1}^\infty (0, 1 - 1/n] = (0, 1)$
is not a finite union of half-open intervals, so $(0,1) \notin \mathcal{A}$.
\end{example}

\begin{definition}[Pre-Measure]
\label{def:premeasure}
A \emph{pre-measure} on an algebra $\mathcal{A}$ is a function
$\mu_0: \mathcal{A} \to [0, \infty]$ satisfying:
\begin{enumerate}[nosep,label=(\roman*)]
    \item $\mu_0(\emptyset) = 0$.
    \item If $A_1, A_2, \ldots \in \mathcal{A}$ are disjoint and
    $\bigsqcup_{n=1}^\infty A_n \in \mathcal{A}$, then
    $\mu_0(\bigsqcup A_n) = \sum \mu_0(A_n)$.
\end{enumerate}
\end{definition}

\subsection{Outer Measures}
\label{subsec:outer}

\begin{definition}[Outer Measure]
\label{def:outer}
An \emph{outer measure} on $\Omega$ is a function
$\mu^*: \powerset(\Omega) \to [0, \infty]$ satisfying:
\begin{enumerate}[nosep,label=(\roman*)]
    \item $\mu^*(\emptyset) = 0$.
    \item \textbf{Monotonicity:} $A \subseteq B \implies \mu^*(A) \leq \mu^*(B)$.
    \item \textbf{Countable sub-additivity:} $\mu^*(\bigcup_{n=1}^\infty A_n) \leq \sum_{n=1}^\infty \mu^*(A_n)$.
\end{enumerate}
\end{definition}

\begin{proposition}[Pre-Measure Induces Outer Measure]
\label{prop:induced_outer}
Given a pre-measure $\mu_0$ on an algebra $\mathcal{A}$, define for any
$E \subseteq \Omega$:
\begin{equation}
    \mu^*(E) = \inf\left\{\sum_{n=1}^\infty \mu_0(A_n) : A_n \in \mathcal{A},\; E \subseteq \bigcup_{n=1}^\infty A_n\right\}.
    \label{eq:induced_outer}
\end{equation}
Then $\mu^*$ is an outer measure, and $\mu^*(A) = \mu_0(A)$ for all $A \in \mathcal{A}$.
\end{proposition}

\begin{proof}
\textit{(i)} $\mu^*(\emptyset) = 0$: cover $\emptyset$ by $A_n = \emptyset$ for all $n$.

\textit{(ii)} If $A \subseteq B$, every cover of $B$ is also a cover of $A$, so the infimum over covers of $A$ is at most the infimum over covers of $B$.

\textit{(iii)} Let $E = \bigcup_n E_n$. For each $\varepsilon > 0$ and each $n$, choose
a cover $\{A_{n,k}\}_{k} \subseteq \mathcal{A}$ with $E_n \subseteq \bigcup_k A_{n,k}$
and $\sum_k \mu_0(A_{n,k}) \leq \mu^*(E_n) + \varepsilon/2^n$. Then
$\{A_{n,k}\}_{n,k}$ covers $E$, so:
\begin{equation}
    \mu^*(E) \leq \sum_n \sum_k \mu_0(A_{n,k}) \leq \sum_n (\mu^*(E_n) + \varepsilon/2^n) = \sum_n \mu^*(E_n) + \varepsilon.
\end{equation}
Since $\varepsilon$ is arbitrary, $\mu^*(E) \leq \sum_n \mu^*(E_n)$.

\textit{$\mu^*$ extends $\mu_0$:} For $A \in \mathcal{A}$, the single-set cover gives
$\mu^*(A) \leq \mu_0(A)$. For the reverse: if $A \subseteq \bigcup_n A_n$ with
$A_n \in \mathcal{A}$, define $B_n = A \cap A_n \cap (\bigcup_{k<n} A_k)^c \in \mathcal{A}$
(since $\mathcal{A}$ is an algebra). Then $A = \bigsqcup B_n$ and $B_n \subseteq A_n$.
By the pre-measure properties: $\mu_0(A) = \sum \mu_0(B_n) \leq \sum \mu_0(A_n)$.
Taking the infimum: $\mu_0(A) \leq \mu^*(A)$.
\end{proof}

\subsection{Carath\'{e}odory's Extension Theorem}
\label{subsec:caratheodory}

\begin{definition}[Carath\'{e}odory Measurability]
\label{def:cara_measurable}
Given an outer measure $\mu^*$ on $\Omega$, a set $A \subseteq \Omega$ is
\emph{$\mu^*$-measurable} (or \emph{Carath\'{e}odory measurable}) if for every
$E \subseteq \Omega$:
\begin{equation}
    \mu^*(E) = \mu^*(E \cap A) + \mu^*(E \cap A^c).
    \label{eq:cara_condition}
\end{equation}
\end{definition}

\begin{remark}
By sub-additivity, $\mu^*(E) \leq \mu^*(E \cap A) + \mu^*(E \cap A^c)$ always holds.
So the content of~\eqref{eq:cara_condition} is the reverse inequality: $A$ ``splits''
every test set $E$ additively. This is a strong condition --- it must hold for
\emph{all} $E \subseteq \Omega$, not just sets in some $\sigma$-algebra.
\end{remark}

\begin{keyresult}[Carath\'{e}odory's Extension Theorem]
\begin{theorem}[Carath\'{e}odory]
\label{thm:caratheodory}
Let $\mu^*$ be an outer measure on $\Omega$. Then:
\begin{enumerate}[nosep,label=(\alph*)]
    \item The collection $\mathcal{M} = \{A \subseteq \Omega : A \text{ is $\mu^*$-measurable}\}$ is a $\sigma$-algebra.
    \item The restriction $\mu = \mu^*|_{\mathcal{M}}$ is a complete measure on $(\Omega, \mathcal{M})$.
\end{enumerate}
Consequently, if $\mu_0$ is a pre-measure on an algebra $\mathcal{A}$ and $\mu^*$ is
the induced outer measure~\eqref{eq:induced_outer}, then $\mathcal{A} \subseteq \mathcal{M}$
and $\mu|_{\mathcal{A}} = \mu_0$. If additionally $\mu_0$ is $\sigma$-finite, then
$\mu$ is the unique extension of $\mu_0$ to $\sigma(\mathcal{A})$.
\end{theorem}
\end{keyresult}

\begin{proof}
\textbf{Part (a): $\mathcal{M}$ is a $\sigma$-algebra.}

\textit{Step 1: $\Omega \in \mathcal{M}$.} For any $E$: $\mu^*(E \cap \Omega) + \mu^*(E \cap \emptyset) = \mu^*(E) + 0 = \mu^*(E)$. \checkmark

\textit{Step 2: Closure under complement.} The condition~\eqref{eq:cara_condition} is symmetric in $A$ and $A^c$, so $A \in \mathcal{M} \Leftrightarrow A^c \in \mathcal{M}$. \checkmark

\textit{Step 3: Closure under finite union.} Let $A, B \in \mathcal{M}$. For any test set $E$:
\begin{align}
    \mu^*(E) &= \mu^*(E \cap A) + \mu^*(E \cap A^c) \tag{$A$ measurable}\\
    &= \mu^*(E \cap A) + \mu^*(E \cap A^c \cap B) + \mu^*(E \cap A^c \cap B^c) \tag{$B$ measurable, applied to $E \cap A^c$}
\end{align}
We need $\mu^*(E) = \mu^*(E \cap (A \cup B)) + \mu^*(E \cap (A \cup B)^c)$.

Note $(A \cup B)^c = A^c \cap B^c$, so the third term is correct. For the first:
$E \cap (A \cup B) = (E \cap A) \cup (E \cap A^c \cap B)$, so by sub-additivity:
$\mu^*(E \cap (A \cup B)) \leq \mu^*(E \cap A) + \mu^*(E \cap A^c \cap B)$.

Combining: $\mu^*(E) \geq \mu^*(E \cap (A \cup B)) + \mu^*(E \cap (A \cup B)^c)$.
The reverse holds by sub-additivity. So $A \cup B \in \mathcal{M}$. \checkmark

\textit{Step 4: Closure under countable disjoint union.} Let $A_1, A_2, \ldots \in \mathcal{M}$ be disjoint. Set $S_n = \bigsqcup_{k=1}^n A_k$ and $S = \bigsqcup_{k=1}^\infty A_k$.

By induction (using Step 3 and closure under complement, hence closure under intersection), $S_n \in \mathcal{M}$, and for any $E$:
\begin{equation}
    \mu^*(E \cap S_n) = \sum_{k=1}^n \mu^*(E \cap A_k).
    \label{eq:cara_finite}
\end{equation}
This is proved by induction: $\mu^*(E \cap S_n) = \mu^*(E \cap S_n \cap A_n) + \mu^*(E \cap S_n \cap A_n^c) = \mu^*(E \cap A_n) + \mu^*(E \cap S_{n-1})$.

Since $S_n \in \mathcal{M}$: $\mu^*(E) = \mu^*(E \cap S_n) + \mu^*(E \cap S_n^c)$.
Since $S_n \subseteq S$, $S^c \subseteq S_n^c$, so $\mu^*(E \cap S_n^c) \geq \mu^*(E \cap S^c)$. Using~\eqref{eq:cara_finite}:
\begin{equation}
    \mu^*(E) \geq \sum_{k=1}^n \mu^*(E \cap A_k) + \mu^*(E \cap S^c).
\end{equation}
Letting $n \to \infty$:
\begin{equation}
    \mu^*(E) \geq \sum_{k=1}^\infty \mu^*(E \cap A_k) + \mu^*(E \cap S^c) \geq \mu^*(E \cap S) + \mu^*(E \cap S^c),
\end{equation}
where the last step uses sub-additivity $\mu^*(E \cap S) \leq \sum \mu^*(E \cap A_k)$.
The reverse inequality holds by sub-additivity. So $S \in \mathcal{M}$. \checkmark

\textit{Step 5: General countable unions.} For arbitrary (not necessarily disjoint)
$A_n \in \mathcal{M}$, write $\bigcup A_n$ as a disjoint union using $B_n = A_n \setminus \bigcup_{k<n} A_k \in \mathcal{M}$, and apply Step 4.

\textbf{Part (b): $\mu$ is a complete measure.}

Countable additivity follows from~\eqref{eq:cara_finite} with $E = S$:
$\mu(S) = \mu^*(S) = \sum \mu^*(A_k) = \sum \mu(A_k)$.

Completeness: if $\mu^*(N) = 0$ and $A \subseteq N$, then for any $E$:
$\mu^*(E \cap A) \leq \mu^*(A) \leq \mu^*(N) = 0$, so
$\mu^*(E \cap A) + \mu^*(E \cap A^c) = \mu^*(E \cap A^c) \leq \mu^*(E)$.
The reverse holds by monotonicity ($E \cap A^c \subseteq E$). Wait, actually we need
$\mu^*(E \cap A^c) \geq \mu^*(E) - \mu^*(E \cap A) = \mu^*(E)$ and $\mu^*(E \cap A^c) \leq \mu^*(E)$. These give equality. So $A \in \mathcal{M}$.

\textbf{Uniqueness under $\sigma$-finiteness:} follows from Corollary~\ref{cor:uniqueness}
applied to the $\pi$-system $\mathcal{A}$, after reducing to the finite-measure case
using the $\sigma$-finite decomposition.
\end{proof}

\subsection{Application: Lebesgue Measure}
\label{subsec:lebesgue}

\begin{theorem}[Existence of Lebesgue Measure]
\label{thm:lebesgue}
There exists a unique measure $\lambda$ on $(\R, \Borel(\R))$ such that
$\lambda((a, b]) = b - a$ for all $a < b$. This measure is $\sigma$-finite and
translation-invariant: $\lambda(A + x) = \lambda(A)$ for all $A \in \Borel(\R)$
and $x \in \R$.
\end{theorem}

\begin{proof}[Proof sketch]
Let $\mathcal{A}$ be the algebra of finite disjoint unions of half-open intervals
$(a_i, b_i]$. Define $\mu_0(\bigsqcup_i (a_i, b_i]) = \sum_i (b_i - a_i)$. One
verifies $\mu_0$ is a pre-measure on $\mathcal{A}$ (the key step is showing
countable additivity, which uses the Heine--Borel theorem: if $(a, b]$ is covered
by countably many $(a_n, b_n]$, then $b - a \leq \sum (b_n - a_n)$).

By Carath\'{e}odory's theorem (Theorem~\ref{thm:caratheodory}), $\mu_0$ extends
uniquely to $\sigma(\mathcal{A}) = \Borel(\R)$ since $\mu_0$ is $\sigma$-finite
($\R = \bigcup_{n=1}^\infty (-n, n]$ with $\mu_0((-n, n]) = 2n < \infty$).
\end{proof}

\begin{theorem}[Non-Measurable Sets (Vitali)]
\label{thm:vitali}
There exists a set $V \subseteq [0, 1]$ that is not Lebesgue measurable.
\end{theorem}

\begin{proof}
Define an equivalence relation on $[0,1)$: $x \sim y$ iff $x - y \in \Q$.
By the Axiom of Choice, select one representative from each equivalence class;
call this set $V$.

For $q \in \Q \cap [0, 1)$, define $V_q = \{x + q \mod 1 : x \in V\}$
(shift $V$ by $q$ modulo 1). The $V_q$ are disjoint (if $v_1 + q_1 \equiv v_2 + q_2$
then $v_1 - v_2 \in \Q$, so $v_1 \sim v_2$, so $v_1 = v_2$ and $q_1 = q_2$).
Also $[0, 1) = \bigsqcup_{q \in \Q \cap [0,1)} V_q$ (every $x$ is equivalent to
some $v \in V$, and $x = v + q$ for some rational $q$).

If $V$ were measurable, then each $V_q$ would have the same measure (by
translation invariance of Lebesgue measure modulo 1). Call this $c$.

By countable additivity: $1 = \lambda([0,1)) = \sum_{q} c$. But $\sum c$ is either
$0$ (if $c = 0$) or $\infty$ (if $c > 0$) --- never $1$. Contradiction.
\end{proof}

\begin{remark}
Vitali's theorem shows why we cannot define Lebesgue measure on $\powerset(\R)$:
not all subsets of $\R$ are measurable. This is the fundamental reason for
restricting to $\sigma$-algebras. In probability, this subtlety is rarely
encountered (finite MDPs use $\powerset(\mathcal{S})$), but it justifies
the formalism and is important for continuous state/action spaces.
\end{remark}

%% ============================================================
%% SECTION 5: MEASURABLE FUNCTIONS
%% ============================================================


\end{document}